\section{Introduction}

Over the past two decades, planetary exploration has entered an era of rapid scientific discovery, marked by unprecedented robotic missions exploring the Moon and Mars. Under the China Lunar Exploration Program (CLEP) and the Tianwen planetary program, the National Astronomical Observatories of the Chinese Academy of Sciences (NAOC) and China National Space Administration (CNSA) have successfully deployed an array of orbital, lander, rover, and sample return vehicles~\cite{jia2021scientific}. These include:
\begin{itemize}
    \item \textbf{Chang'e-1 (2007) and Chang'e-2 (2010)}: Global lunar orbiters providing high-resolution stereoscopic imagery (down to 7~m globally and 1.3~m locally), laser altimetry topography (LAM), microwave radiometer brightness temperature maps (CELMS), and asteroid (4179) Toutatis close flyby imaging.
    \item \textbf{Chang'e-3 (2013)}: The first Chinese lunar soft landing in Mare Imbrium (Zi Wei Crater), deploying the \textit{Yutu} rover equipped with the Lunar Penetrating Radar (LPR) and Visible and Near-Infrared Spectrometer (VNIS)~\cite{su2014lunar}.
    \item \textbf{Chang'e-4 (2019)}: The historic first soft landing on the lunar farside within Von K\'{a}rm\'{a}n Crater in the South Pole-Aitken (SPA) basin, operating the \textit{Yutu-2} rover for over five years to probe deep lunar crustal and mantle materials~\cite{lai2020first,li2020change,wimmer2020first}.
    \item \textbf{Chang'e-5 (2020) and Chang'e-6 (2024)}: Groundbreaking sample return missions retrieving 1,731~g of young ($~2.0$~Ga) basalts from Mons R\"{u}mker (Oceanus Procellarum) and 1,935.3~g of ancient farside basalts from Apollo Basin, demonstrating in-situ detection of lunar hydroxyl ($OH/\text{H}_2\text{O}$) by the Lunar Mineralogical Spectrometer (LMS)~\cite{li2022situ}.
    \item \textbf{Tianwen-1 and Zhurong Rover (2021)}: China's maiden Mars mission, deploying the \textit{Zhurong} rover in Southern Utopia Planitia to investigate subsurface sedimentary flood layers (RoRAD), laser-induced breakdown spectroscopy of hydrated minerals (MarSCoDe), surface crustal magnetism (RoMAG), and Martian boundary layer meteorology (MCS)~\cite{zhang2022layered,liu2022zhurong,peng2022mars}.
\end{itemize}

The raw and calibrated observational data generated by these missions are formally archived in the \textbf{Lunar and Planetary Data Release System (CLPDS)} (\url{https://clpds.bao.ac.cn}). However, accessing, synthesizing, and interpreting these vast multi-instrument datasets presents substantial technical challenges. Data products are distributed across varying processing levels (PDS Level 0 to Level 4) in specialized binary and tabular structures (such as .2BL, .2C, .IMG, and .DAT formats), often requiring specialized proprietary calibration pipelines.

To address this gap and promote FAIR data principles (Findable, Accessible, Interoperable, and Reusable), we developed the \textbf{CLPDS Planetary Exploration Data Visualization Suite}. This open-source framework couples scientific data ingestion with modern client-side WebGL/Three.js computing, enabling interactive exploration of planetary surfaces, subsurface stratigraphy, in-situ spectroscopy, and planetary environmental dynamics.
